\documentclass{article}

\usepackage{amsmath,amssymb}
\usepackage{xurl}
\usepackage[hidelinks]{hyperref}
\setlength{\emergencystretch}{2em}

% Shared commands for notes in docs/paper-gaps/.

\DeclareUnicodeCharacter{2080}{\ifmmode{}_{0}\else\textsubscript{0}\fi}
\DeclareUnicodeCharacter{2081}{\ifmmode{}_{1}\else\textsubscript{1}\fi}
\DeclareUnicodeCharacter{2082}{\ifmmode{}_{2}\else\textsubscript{2}\fi}
\DeclareUnicodeCharacter{2099}{\ifmmode{}_{n}\else\textsubscript{n}\fi}

\newcommand{\Lean}{\textsc{Lean}}
\ExplSyntaxOn
\NewDocumentCommand{\leanid}{m}{
  \ifmmode
    \textsf{\detokenize{#1}}
  \else
    \group_begin:
    \urlstyle{sf}
    \tl_set:Nn \l_tmpa_tl {#1}
    \tl_replace_all:Nnn \l_tmpa_tl {\_} {_}
    \exp_args:NV \path \l_tmpa_tl
    \group_end:
  \fi
}
\ExplSyntaxOff
\providecommand{\tr}{\operatorname{tr}}
\newcommand{\tnleanIssueBase}{https://github.com/LionSR/TNLean/issues/}
\newcommand{\tnleanPrBase}{https://github.com/LionSR/TNLean/pull/}
\newcommand{\tnleanDocBase}{https://sirui-lu.com/TNLean/docs/}
\newcommand{\ghissue}[1]{\href{\tnleanIssueBase#1}{GitHub issue \##1}}
\newcommand{\ghpr}[1]{\href{\tnleanPrBase#1}{PR \##1}}
\newcommand{\doclink}[1]{\href{\tnleanDocBase#1.html}{\path{#1}}}

% Dirac notation and the identity operator, as in blueprint/src/macros/common.tex.
% \providecommand keeps note-local definitions authoritative.
\usepackage{dsfont}
\providecommand{\ket}[1]{|#1\rangle}
\providecommand{\bra}[1]{\langle #1|}
\providecommand{\braket}[2]{\langle #1 | #2 \rangle}
\providecommand{\ketbra}[2]{|#1\rangle\!\langle #2|}
\providecommand{\Id}{\mathds{1}}
\providecommand{\one}{\mathds{1}}
\providecommand{\C}{\mathbb{C}}
\providecommand{\MN}[1]{M_{#1}(\C)}
\providecommand{\MD}{M_D(\C)}

% Verdict marker: \gapnote{<kind>}{<status>}. Kind is the policy
% classification (clarification, local-correction, scope-restriction,
% unfaithful, false-source, open-gap); status is open, wip (elimination
% underway), resolved, or historical. Machine-read by the site index;
% prints a verdict line.
\newcommand{\gapnote}[2]{\par\noindent\textbf{Verdict:} #1 (#2).\par}


% Fallbacks keep this note compilable when a different command.tex is found first.
\providecommand{\leanid}[1]{\textsf{\detokenize{#1}}}
\providecommand{\gapnote}[2]{\par\noindent\textbf{Verdict:} #1 (#2).\par}
\providecommand{\ket}[1]{|#1\rangle}
\providecommand{\bra}[1]{\langle #1|}

\title{Party Ownership in the Elimination of Normalized Pair Effects}
\author{}
\date{2026-10-07}

\begin{document}
\maketitle
\gapnote{scope-restriction}{open}

\section{The source lemma}

Lemma~5.1 of the polynomial PEPS approximation manuscript of September~24,
2026 (\path{04-compression.tex}, lines~53--71) considers a contraction
$G=\sum_\xi c_\xi M_\xi$ expanded into $K$ allowed monomials, each with at
most $r$ pair effects. An allowed monomial (lines~32--35) is a composition of
local contractions, preparations of normalized vectors shared between two
participating parties, and contractions by normalized bras shared between two
participating parties. For every $m\ge1$ the lemma produces a gate
$\widetilde G_m$ and a normalized vector $\Gamma_m$ with
\[
  \bigl\|\widetilde G_m-G\otimes\ket{\Gamma_m}\bigr\|
  \le\frac{r}{\sqrt m}\sum_\xi|c_\xi|,
\]
and asserts four further properties of $\widetilde G_m$: (i) it has an
expansion using only local contractions and normalized pair sources; (ii) the
expansion has at most $Km^r$ monomials with absolute coefficient sum at most
$\sum_\xi|c_\xi|$; (iii) its additional registers are owned by the original
participating parties; (iv) all sources on one pair of parties may be combined
into one normalized pair source.

\section{The model without parties}

The formal model writes a monomial with $n$ effects as a chain
$F_n\circ E_n\circ\cdots\circ E_1\circ F_0$, where each $E_i=\bra{\eta_i}\otimes\mathbf{1}$
is a recorded normalized pair effect and each $F_i$ is an arbitrary
contraction. The model has no parties. The error bound, the normalization of
$\Gamma_m$, the common additional output space of all branches, the count
$Km^r$, the equality of the absolute coefficient sums, and the contractivity of
every term of the expansion are proved.\footnote{The packaged statement is
\leanid{TNLean.PEPS.PairEffect.pairEffectElimination}; the rescaling at a
chosen stack length is
\leanid{TNLean.PEPS.PairEffect.replaceGate_rescaled_stackLength}.} For (iv),
the formal statement is the norm identity
$\|\eta\otimes\eta'\|=\|\eta\|\,\|\eta'\|$ for pair vectors regrouped on the
tensor products of their half-spaces.

In that model clauses (i) and (iii) cannot be stated. Without parties, a
local contraction cannot be told apart from an arbitrary contraction, and a
pair effect is itself a contraction; a pair effect may therefore sit inside
some $F_i$, and a predicate ``source-only'' formulated without parties would
hold of every contraction.

\section{Discharge}

The monomials are now placed on a party layout. A register has an owner party
and a Hilbert space, and the memory is the tensor product of the registers.
A source-only word is a composition of local contractions, each acting on
registers of a single party next to arbitrary other registers, of
preparations of normalized vectors on two fresh registers of two distinct
parties, and of reorderings of tensor factors. An allowed monomial is a chain
of source-only words separated by pair effects, each effect contracting a
register $\mathbb C^\alpha$ of a party $p$ and a register $\mathbb C^\beta$ of
a party $q\ne p$ with a normalized bra. Its operator is a chain of the form
above, so the estimates already proved apply to it. Writing the effect on registers
$\mathbb C^\alpha$ and $\mathbb C^\beta$ at the front of the list is a
convention: an effect on finite-dimensional registers $\mathcal U$ of $p$ and
$\mathcal V$ of $q$ in any position takes this form after local isometric
equivalences $\mathcal U\cong\mathbb C^\alpha$ and
$\mathcal V\cong\mathbb C^\beta$ and a reordering of tensor factors, all
absorbed into the preceding source-only word without changing its norm or the
coefficient of the monomial.

In the replaced monomial, the stack of an effect occurrence on $(p,q)$
consists of $m$ pair registers, the $\mathbb C^\alpha$ half of each owned by
$p$ and the $\mathbb C^\beta$ half owned by $q$. The insertion of the input at
position $k$ among $m-1$ copies of $\eta$ is a source-only word: pair sources
of $\eta$ are prepared in the other positions, so no register permutation is
needed. Under the identification of the stack with
$(\mathbb C^{\alpha\times\beta})^{\otimes m}$ this word is the insertion map
$I_k$. Consequently:
\begin{itemize}
\item every term of the expansion of $\widetilde G_m$, read on the output
  layout, is the operator of a source-only word with normalized sources and
  contractive local maps (clause (i));
\item the output layout is the list of the stacks of all effect occurrences,
  followed by the output layout of $G$, and the stack of an effect on
  $(p,q)$ consists of $m$ registers of $p$ alternating with $m$ registers of
  $q$ (clause (iii));
\item the garbage vector $\Gamma_m$ is prepared by normalized pair sources on
  the endpoint parties;
\item two pair sources on the same pair of parties, prepared one after the
  other and followed by merging the two registers of each party, are the
  single pair source of the regrouped vector, whose norm is the product of the
  norms.
\end{itemize}
The packaged statement of clauses (i)--(iii) and the bound is
\leanid{TNLean.PEPS.PairEffect.partyPairEffectElimination}. Its hypotheses
are those of the source lemma, without the contraction hypothesis on $G$,
which the statement does not use. The participating parties of clause (iii)
are read as the endpoint parties of the effects of the monomials. The
combination of two adjacent sources is
\leanid{TNLean.PEPS.PairEffect.eval_combineSources}.

\section{Remaining restriction: combining sources}

Clause (iv) is applied to the expansion of $\widetilde G_m$, whose words
contain about $m$ sources for every effect occurrence, interleaved with local
contractions; Theorem~5.2 (lines~132--212) needs the number of pair sources
per monomial bounded independently of $m$. The source's proof of clause (iv)
(lines~125--127) is the tensor-product identity above. The formal statement
covers only two sources prepared one after the other. The step that every
pair source prepares fresh registers and can therefore be moved past the
operations on other registers, so that all sources on one pair of parties
become adjacent, is not formalized.

Elimination: prove that every allowed source-only word has the operator of
the preparation of all its sources followed by a word without sources, then
reorder the prepared registers and merge the sources on each pair of parties
one at a time with the identity above. The inverse merging maps are local
isometries, so the resulting word stays allowed.

\end{document}
